% !TeX spellcheck = it_IT

\documentclass[handout]{beamer}

\usepackage[utf8]{inputenc}
\usepackage[italian]{babel}
\usepackage{multicol}
\usetheme{metropolis}

\usepackage{listings}

\usepackage{tikz}
\usetikzlibrary{shapes,arrows,positioning}

\tikzstyle{startstop} = [ellipse, draw, fill=bg, text=structure, text width=3em, text badly centered, inner sep=0pt]
\tikzstyle{block} = [rectangle, draw, fill=bg, text=structure, text width=5em, text centered]
\tikzstyle{decision} = [diamond, aspect=2, draw, fill=bg, text=structure, text width=3em, text badly centered, inner sep=0pt]
\tikzstyle{input} = [rectangle, draw, fill=bg, text=structure, text width=4em, text centered, rounded rectangle, rounded rectangle east arc=none]
\tikzstyle{output} = [rectangle, draw, fill=bg, text=structure, text width=4em, text centered, rounded rectangle, rounded rectangle west arc=none]

\tikzstyle{line} = [draw, -latex']


\title{Informatica B - Esercitazione 3}
\subtitle{Array}

\author{Stefano Cereda\\
		stefano1.cereda@mail.polimi.it
	}
\date{17/10/2017}
\institute[PoliMi]{\vspace{0.5cm}\centering Politecnico di Milano \\ \vspace{0.2cm}
	\includegraphics[width=0.2\linewidth]{./logopolimi}}
%\logo{\includegraphics[width=15mm]{./logopolimi}}

\setbeamercovered{invisible}

\makeindex

\begin{document}
\begin{frame}
	\maketitle
\end{frame}


\begin{frame}{Esempio 1 - da slide lezione}
Scrivere un programma che legga un vettore di interi $V$ ed un numero intero $n$, si stampi l'indice di $V$ corrispondente al primo elemento uguale ad $n$.
\end{frame}

\begin{frame}{Esercizio 2}
Dato un vettore ordinato di interi $V$ ed un numero $n$, si scriva la porzione di codice per implementare un algoritmo di ricerca e stampare l'indice del primo elemento di $V$ uguale ad $n$. (Esempio 1)

Si modifichi il programma per trovare più occorrenze dello stesso numero (stampandone tutti gli indici).

Si modifichi il programma per creare un secondo array $V_2$ in cui inserire gli indici delle occorrenze (quelli stampati al punto precedente). Si presti attenzione a non lasciare buchi in $V_2$. Al termine si stampi $V_2$.
\end{frame}

\begin{frame}{La libreria string.h}
La libreria string.h vi offre alcune funzioni per eseguire operazioni comuni sulle stringhe. Ad esempio:
\begin{itemize}
	\item Ottenere la lunghezza di una stringa
	\item Confrontare due stringhe
	\item Concatenare due stringhe
	\item Cercare un carattere in una stringa
	\item Cercare una stringa in un'altra
	\item Copiare una stringa in un'altra
\end{itemize}
\end{frame}

\begin{frame}{La libreria string.h - Vedere il codice per gli esempi}
\begin{itemize}
	\item \texttt{int strlen(char* str);} $\rightarrow$ numero di caratteri in \texttt{str} senza contare il terminatore
	\item \texttt{int strcmp(char* str1, char* str2);} $\rightarrow$ $<$0 se \texttt{str1} precede lessicograficamente \texttt{str2}, $>$0 per il viceversa, =0 se sono uguali
	\item \texttt{char* strcat(char* dest, char* src);} $\rightarrow$ appende \texttt{str} in coda a \texttt{dest}
	\item \texttt{char* strchr(char* str, int car);} $\rightarrow$ \texttt{car} all'interno di \texttt{str} e restituisce la porzione di \texttt{str} successiva alla prima occorrenza di \texttt{car} (incluso) oppure \texttt{NULL}
	\item \texttt{char* strstr(char* str, char* sub);} $\rightarrow$ \texttt{sub} all'interno di \texttt{str} e restituisce la porzione di \texttt{str} successiva alla prima occorrenza di \texttt{sub} (inclusa) oppure \texttt{NULL}
	\item \texttt{char* strcpy(char* dest, char* src)} $\rightarrow$ copia \texttt{src} in \texttt{dest e restituisce \texttt{dest}}
\end{itemize}
\end{frame}

\begin{frame}{Esercizio 3}
Scrivere un programma che riceva un numero noto $N=R\cdot C$ di interi e li salvi in una matrice $M$ di dimensioni note $R \times C$. Si stampi poi la matrice in formato tabellare.
\end{frame}

\begin{frame}{Esercizio 4}
Data una matrice quadrata intera $M$, di dimensioni $D \times D$, si scriva il codice per calcolare e visualizzare la trasposta. (le righe diventano colonne e viceversa).
\end{frame}

\begin{frame}{Esercizio 5}
Date due matrici intere $A, B$ si scriva un programma che ne calcoli il prodotto $C=A\times B$

Guardare le slide della lezione per la definizione. Praticamente ricordatevi ``mano sinistra in orizzontale si muove verso destra e mano destra in verticale si muove verso il basso''.
\end{frame}

\end{document}